
\section{Introduction}
\label{sec:introduction}

Large language models (LLMs) are increasingly used in systems that
generate text, assist with decisions, and interact with digital
environments. Their usefulness also creates safety challenges: a model
may produce harmful instructions, conceal its behavior, manipulate an
evaluation process, or otherwise act contrary to the system's intended
objectives. Reliable monitoring is therefore an important part of
evaluating and deploying LLM-based systems.

A monitor can inspect different parts of a model interaction. An
output-only monitor receives the user's prompt and the model's final
answer. A more information-rich monitor may also receive an available
reasoning trace, often referred to as chain-of-thought (CoT). In
principle, additional context may expose signals that are not visible in
the final answer alone. However, access to a reasoning trace does not
automatically guarantee more accurate detection: the trace may be
incomplete, noisy, unavailable, or may not faithfully represent the
model's internal computation.

This creates the research problem addressed in this project: whether
including a reasoning trace changes the ability of a monitor to identify
unsafe behavior under a controlled evaluation setup. We compare two
conditions: (1) output-only monitoring, which uses the prompt and final
output, and (2) CoT-plus-output monitoring, which uses the prompt,
reasoning trace, and final output. The study focuses on four categories
of unsafe behavior: reward hacking, deception, sabotage, and harmful
compliance.

The research gap is the need to evaluate the effect of the information
available to a monitor while keeping the monitoring approach and
evaluation conditions explicit. This project provides a controlled
comparison of the two input conditions using the benchmark data and
keyword-based monitoring implementation available in the repository.
Because the current monitor is a keyword baseline, the study should be
interpreted as an evaluation of this baseline and dataset, not as proof
that reasoning traces universally improve LLM safety monitoring.

The research questions are:

\begin{itemize}
    \item \textbf{RQ1:} How does access to a reasoning trace affect the
    monitor's ability to detect unsafe behavior compared with
    output-only monitoring?

    \item \textbf{RQ2:} What is the recall of each monitoring condition
    at a target false-positive rate of 5\%?

    \item \textbf{RQ3:} How do the monitoring conditions compare in
    operational performance, particularly latency, under the tested
    request loads?
\end{itemize}

The contributions of this project are:

\begin{itemize}
    \item A comparison of output-only and CoT-plus-output monitoring
    conditions.

    \item An evaluation framework focused on unsafe-behavior categories
    including reward hacking, deception, sabotage, and harmful compliance.

    \item An analysis plan that considers detection quality at a fixed
    false-positive rate and operational latency.

    \item A reproducible experimental configuration and processing
    workflow for the benchmark.
\end{itemize}

The remainder of the paper describes the experimental methodology,
reports the measured results, discusses their implications and
limitations, and concludes with directions for future work.
